\subsection{PROJECTIVE MODULES}

DEFINITION 1. An A-module P is called projective if, for every exact sequence \(\mathbf{F}' \to \mathbf{F} \to \mathbf{F}''\) of A-linear mapping, the sequence

\[
\operatorname{Hom} (\mathrm{P}, \mathrm{F} ^ {\prime}) \rightarrow \operatorname{Hom} (\mathrm{P}, \mathrm{F}) \rightarrow \operatorname{Hom} (\mathrm{P}, \mathrm{F} ^ {\prime \prime})
\]

is exact.

PROPOSITION 3. For an A-module P, the direct sum of a family of submodules \((\mathbf{M}_{t})\) , to be projective, it is necessary and sufficient that each of the \(\mathbf{M}_{t}\) be projective.

For every A-module homomorphism \(u: \mathbf{E} \to \mathbf{F}\) ,

\[
\operatorname{Hom} \left(1 _ {\mathrm{P}}, u\right): \operatorname{Hom} (\mathrm{P}, \mathrm{E}) \rightarrow \operatorname{Hom} (\mathrm{P}, \mathrm{F})
\]

is identified with \(\prod_{i}\operatorname{Hom}(1_{M_i},u)\) (§ 1, no. 6, Corollary 1 to Proposition 6); the conclusion thus follows from Definition 1 and § 1, no. 5, Proposition 5 (ii).

COROLLARY. Every free A-module is projective.

It suffices by Proposition 3 to show that \(A_{s}\) is projective, which follows immediately from the commutativity of diagram (50) of §1, no. 14.

PROPOSITION 4. Let P be an A-module. The following properties are equivalent:

\begin{enumerate}
\def\labelenumi{(\alph{enumi})}
\item
  P is projective.
\item
  For every exact sequence \(0 \to \mathbf{F}' \to \mathbf{F} \to \mathbf{F}'' \to 0\) of A-linear mappings, the sequence
\end{enumerate}

\[
0 \rightarrow \operatorname{Hom} (\mathrm{P}, \mathrm{F} ^ {\prime}) \rightarrow \operatorname{Hom} (\mathrm{P}, \mathrm{F}) \rightarrow \operatorname{Hom} (\mathrm{P}, \mathrm{F} ^ {\prime \prime}) \rightarrow 0
\]

is exact.

\begin{enumerate}
\def\labelenumi{(\alph{enumi})}
\setcounter{enumi}{2}
\item
  For every surjective A-module homomorphism \(u: \mathbf{E} \to \mathbf{E}''\) and every homomorphism \(f: \mathbf{P} \to \mathbf{E}''\) , there exists a homomorphism \(g: \mathbf{P} \to \mathbf{E}\) such that \(f = u \circ g\) (it is said that \(f\) can be ``lifted'' to a homomorphism of \(\mathbf{P}\) into \(\mathbf{E}\) ).
\item
  Every exact sequence \(0 \to \mathbf{E}' \to \mathbf{E} \xrightarrow{v} \mathbf{P} \to 0\) of A-linear mappings splits (and therefore \(\mathbf{P}\) is isomorphic to a direct factor of \(\mathbf{E}\) ).
\item
  P is isomorphic to a direct factor of a free A-module.
\end{enumerate}

It is trivial that (a) implies (b). To see that (b) implies (c), it suffices to apply (b) to the exact sequence \(0 \to \mathbf{E}' \to \mathbf{E} \xrightarrow{u} \mathbf{E}'' \to 0\) , where \(\mathbf{E}' = \operatorname{Ker}(u)\) , since (c) expresses the fact that

\[
\operatorname{Hom} \left(\mathbf {1} _ {\mathrm{P}}, u\right): \operatorname{Hom} (\mathrm{P}, \mathrm{E}) \rightarrow \operatorname{Hom} (\mathrm{P}, \mathrm{E} ^ {\prime \prime})
\]

is surjective. To see that (c) implies (d), it suffices to apply (c) to the surjective homomorphism \(v: \mathbf{E} \to \mathbf{P}\) and the homomorphism \(\mathbf{l}_{\mathbf{P}}: \mathbf{P} \to \mathbf{P}\) ; the existence of a homomorphism \(g: \mathbf{P} \to \mathbf{E}\) such that \(\mathbf{l}_{\mathbf{P}} = v \circ g\) implies that the sequence

\[
0 \longrightarrow \mathrm{E} ^ {\prime} \longrightarrow \mathrm{E} \stackrel {v} {\longrightarrow} \mathrm{P} \longrightarrow 0
\]

splits (§ 1, no. 9, Proposition 15). As for every A-module M there exist a free A-module L and an exact sequence \(0 \to R \to L \to M \to 0\) (§ 1, no. 11, Proposition 20), clearly (d) implies (e). Finally (e) implies (a) by virtue of Proposition 3 and its Corollary.

COROLLARY 1. For an A-module to be projective and finitely generated, it is necessary and sufficient that it be a direct factor of a free A-module with a finite basis.

The condition is obviously sufficient; conversely, a finitely generated projective module E is isomorphic to a quotient of a free module F with a finite basis (§ 1, no. 11) and E is isomorphic to a direct factor of F by virtue of Proposition 4 (d).

COROLLARY 2. Let C be a commutative ring and E, F two finitely generated projective C-modules; then \(\operatorname{Hom}_{\mathbf{C}}(\mathbf{E}, \mathbf{F})\) is a finitely generated projective C-module.

It may be assumed that there are two finitely generated free C-modules M, N such that \(M = E \oplus E'\) , \(N = F \oplus F'\) ; it follows from §1, no. 6, Corollary 1 to Proposition 6 that \(\operatorname{Hom}_{\mathbb{C}}(\mathbf{M}, \mathbf{N})\) is finitely generated and free and on the other hand that \(\operatorname{Hom}_{\mathbb{C}}(\mathbf{M}, \mathbf{N})\) is isomorphic to

\[
\operatorname{Hom} _ {\mathbf {c}} (\mathbf {E}, \mathbf {F}) \oplus \operatorname{Hom} _ {\mathbf {c}} (\mathbf {E} ^ {\prime}, \mathbf {F}) \oplus \operatorname{Hom} _ {\mathbf {c}} (\mathbf {E}, \mathbf {F} ^ {\prime}) \oplus \operatorname{Hom} (\mathbf {E} ^ {\prime}, \mathbf {F} ^ {\prime}),
\]

whence the corollary.

